\section{(Co)lax localizations}
\label{sec:laxloc}

\gutentag{00C7}%
We will first recall Bousfield localizations, whose monad is idempotent
(\cref{subsec:laxloc-bousfield}). In this section the monad is lax idempotent
(\cref{subsec:prelim-monads}): we study \emph{lax localizations}. We will show
that a lax localization makes a class of $1$-morphisms coreflective and a class
of squares between them right Beck--Chevalley (\cref{prop:laxloc-class}), that
its modules are the local objects for these classes, and that the modules and
their $1$-morphisms are generated by the free ones under two closure rules for
reflective $1$-morphisms (\cref{thm:laxloc-local}). Conversely,
in the presence of
colax limits of $1$-morphisms, a left adjoint to the inclusion of the local
objects for any class is a lax localization (\cref{cor:laxloc-P}), which will
give a bijection with the classes that equal their saturation
(\cref{cor:laxloc-classes}). We will also show that the right adjoint induces
coreflective functors on the hom-categories out of objects in the image of the
left adjoint (\cref{prop:laxloc-homs}).
Colax localizations are dual (\cref{rmk:laxloc-colax}). The two-dimensional
localizations of \cref{sec:2dloc} will make one class reflective and another
coreflective at the same time.

\subsection{Recollection: Bousfield localization}
\label{subsec:laxloc-bousfield}

\gutentag{0028}%
We recall the one-dimensional situation that we categorify.

\begin{defi}
\gutentag{0029}%
A \emph{Bousfield localization} of an $ \infty $-category $ C $
is a localization $ L: C \to D $ in the $ (\infty,
2) $-category of $ \infty $-categories.
\end{defi}

\gutentag{002A}%
A Bousfield localization can be described as a category obtained from $ C $ by
formally
inverting a class of morphisms.

\subsubsection*{Local objects}
\gutentag{002B}%
Let $\C$ be an $\infty$-category and $S$ a class of morphisms of $\C$.
An object $c$ is \emph{$S$-local} if
$s^{\natural}\colon \mapp_{\C}(y,c) \to \mapp_{\C}(x,c)$
is an equivalence for every
$(s\colon x \to y) \in S$;
let
$\iota\colon \C_{S} \subset \C$
be the full subcategory of $S$-local objects.
By the Yoneda lemma, $s$ gives a functor
$R_{s}\colon \C \to \func(\Delta^{1},\Spc)$,
sending $ c $ to
$s^{\natural}\colon \mapp_{\C}(y,c) \Rightarrow \mapp_{\C}(x,c)$.
The inclusion of the equivalences
$\Spc \subset \func(\Delta^{1},\Spc)$
is fully faithful, and pullbacks of fully faithful functors are fully faithful,
so
\begin{equation}
\label{eq:bousfield-pullback}
\C_{S} \simeq \C \times_{\prod_{S}\func(\Delta^{1},\Spc)}
\textstyle\prod_{S}\Spc \ .
\end{equation}
Its two-dimensional form will be \cref{thm:local-pullback}.
Here the walking equivalence is contractible,
so being an equivalence is a property and there is no condition on morphisms.

\subsubsection*{Localizations}
\gutentag{00A2}%
Let $W$ be the class of \emph{$S$-local equivalences},
the $f$ such that every $S$-local object is $\{f\}$-local.
Then $S \subset W$ and $\C_{W} = \C_{S}$. As $W$ is large in general,
we never form a pushout along $W$; the localization at $W$ is characterized by
its universal property instead.

\gutentag{002C}%
Let
$\Edge_{1} := \func(\Delta^{1},\Spc) \simeq \Free(\Delta^{1})$
(presheaves of spaces on
$\Delta^{1} \simeq (\Delta^{1})\opcat$)
be the free presentable $\infty$-category on an arrow, so
$\func^{L}(\Edge_{1},\C) \simeq \func(\Delta^{1},\C)$
for $\C \in \prl$. The free presentable $\infty$-category on an
\emph{equivalence} is $\Eqc := \Spc$.
The functor
$\Delta^{1} \to \mathrm{pt}$
induces
$e_{!}\colon \Edge_{1} \to \Eqc$,
$(A \to B) \mapsto B$, with fully faithful right adjoint
$e^{*}\colon B \mapsto (B = B)$.

\begin{prop}
\label{prop:bousfield-pushout}
\gutentag{002E}%
Let $\C$ be an $\infty$-category and $S$ a class of morphisms of $\C$.
\begin{enumerate}
\item If
$\iota\colon \C_{S} \subset \C$
has a left adjoint $L$, then for every $\infty$-category $\E$,
restriction along $L$ identifies $\func(\C_{S},\E)$ with the full subcategory of
$\func(\C,\E)$ of functors inverting $W$ \cite[5.2.7.12]{htt}:
$L$ is the localization of $\C$ at $W$. Moreover $T := \iota L$ is an idempotent
monad and
$\C_{S} \simeq \modcat{T}{\C}$.
\item If $\C$ is presentable and $S$ is a set, then $\iota$ has a left adjoint,
$W$ is the strongly saturated class generated by $S$,
and $\C_{S}$ is presentable \cite[5.5.4.15]{htt}. Viewing $S$ as a map
$\bigoplus_{S}\Edge_{1} \to \C$
in $\prl$ ($\bigoplus$ the coproduct in $\prl$), $\C_{S}$ is the pushout
\begin{equation*}
\C \sqcup_{\bigoplus_{S}\Edge_{1}} \textstyle\bigoplus_{S}\Eqc
\end{equation*}
in $\prl$: $L$ is universal among colimit-preserving functors inverting $S$.
\end{enumerate}
\end{prop}

\begin{proof}
\gutentag{002F}%
(1) By the Yoneda lemma, $Lf$ is invertible if and only if $f \in W$,
and $\iota$ is fully faithful, so \cite[5.2.7.12]{htt} applies.
As $\iota$ is fully faithful, the counit
$L\iota \Rightarrow \mathrm{id}$
and hence
$\mu = \iota\varepsilon L$
are equivalences; a fully faithful right adjoint is monadic.
(2) The first claims are \cite[5.5.4.15]{htt}. For $\D \in \prl$,
$\func^{L}(\Edge_{1},\D) \simeq \func(\Delta^{1},\D)$
and
$\func^{L}(\Eqc,\D) \simeq \D$
is the full subcategory of equivalences,
so the pushout has the stated universal property. Under
$\prl \simeq (\prr)\opcat$
\cite[5.5.3.4]{htt} it becomes a pullback in $\prr$,
computed in $\widehat{\Cat}$ \cite[5.5.3.18]{htt}. The right adjoint of
$s\colon \Edge_{1} \to \C$
is $R_{s}$ and that of $e_{!}$ is $e^{*}$,
so this pullback is \eqref{eq:bousfield-pullback}.
\end{proof}

\gutentag{00A3}%
In (1), $W$ cannot be replaced by $S$, and in (2) the pushout in $\prl$ cannot
be replaced by the pushout
$\C \sqcup_{\coprod_{S}\Delta^{1}} \coprod_{S}\mathrm{pt}$
in $\Cat$. For the poset
$\C = \{0 < a < t,\ 0 < b < t\}$
and $S = \{0 \to a\}$ we have $\C_{S} = \{a < t\}$,
while the pushout in $\Cat$ is $\{a < b < t\}$: it does not invert $b \to t$,
the pushout of $0 \to a$ along $0 \to b$,
because functors out of it need not preserve pushouts.
So (2) needs colimit-preserving functors, and (1) needs the whole class $W$.

\gutentag{0030}%
Presentability is used only for the existence of $L$ and for (2).
The categorification will follow the same split.
It will replace the walking equivalence (contractible) by the walking
coreflection of \cref{subsec:mates-morphisms} in this section, and by the
walking reflection and coreflection in \cref{sec:2dloc} (not contractible,
but being reflective or coreflective is still a property by
\cref{lem:mates-refl}). As inverses, unlike adjoints,
are compatible with every commutative square,
the data will also include squares: $P$, and in \cref{sec:2dloc} also $I$, will
be locally full sub-$(\infty,2)$-categories of $\func(\Delta^{1},\D)$,
whose $1$-morphisms are to become Beck--Chevalley squares.
Local objects will be defined and studied in an arbitrary $(\infty,2)$-category
in \cref{subsec:laxloc-local,subsec:2dloc-local,subsec:2dloc-plain},
saturated classes in \cref{subsec:2dloc-saturation},
and presentability will be used for the existence of the left adjoint in
\cref{subsec:2dloc-presentable}.
The local objects will form only a \emph{locally full} subcategory,
and the monad will not be idempotent. Unlike in (1),
the local objects will in general not be the formal localization in $\twoCat$,
even when the left adjoint exists; \cref{prop:local-vs-free} will describe the
difference, and \cref{ex:local-vs-pushout} will exhibit it.

\subsection{Lax localizations}
\label{subsec:laxloc-def}

\gutentag{00C8}%
An \emph{adjunction of $(\infty,2)$-categories} is an adjunction
(\cref{def:adjunction}) in $\twoCat$, regarded as an $(\infty,2)$-category with
hom-categories
$\iota_{1}\func(\A,\B)$:
the non-invertible $2$-morphisms of $\func(\A,\B)$ are discarded, not inverted.

\gutentag{00C9}%
Throughout the rest of this section, $L \dashv R$ is an adjunction of
$(\infty,2)$-categories, with
$L\colon \D \to \C$,
$R\colon \C \to \D$
and unit
$\eta\colon \mathrm{id}_{\D} \Rightarrow RL$.
We write
$\varepsilon\colon LR \Rightarrow \mathrm{id}_{\C}$
for its counit and
\begin{equation*}
\varphi\colon \varepsilon L \circ L\eta \simeq \mathrm{id}_{L} \ , \qquad
\theta\colon R\varepsilon \circ \eta R \simeq \mathrm{id}_{R}
\end{equation*}
for the invertible $2$-morphisms of $\func(\D,\C)$ and $\func(\C,\D)$ given by
the triangle identities.
For $d \in \D$ and $c \in \C$ the composite
\begin{equation}
\label{eq:laxloc-adjunction}
\Hom_{\C}(Ld,c) \xrightarrow{\ R\ } \Hom_{\D}(RLd,Rc)
\xrightarrow{\ \eta_{d}^{\natural}\ } \Hom_{\D}(d,Rc)
\end{equation}
is an equivalence, natural in $d$ and $c$ as functors
$\D\opcat \times \C \to \Cat$:
an inverse is
$g \mapsto \varepsilon_{c} \circ Lg$,
by naturality of $\eta$ and
$\varepsilon$ and by $\varphi$ and $\theta$.
Restricting the adjunction
$\Adj \to \twoCat$
to the full sub-$(\infty,2)$-category of $\Adj$ on $-$, the walking monad
\cite[Not.~5.1]{haugseng2021laxmonads}, gives a monad $T := RL$ on $\D$
(\cref{def:monad}) by \cite[Cor.~8.9]{haugseng2021laxmonads}, with unit
$\eta$, multiplication
$\mu := R\varepsilon L$
and unit laws
\begin{equation*}
\theta L\colon \mu \circ \eta T \simeq \mathrm{id}_{T} \ , \qquad
R\varphi\colon \mu \circ T\eta \simeq \mathrm{id}_{T} \ .
\end{equation*}

\begin{defi}
\label{def:laxloc}
\gutentag{00CA}%
The adjunction $L \dashv R$ is a \emph{lax localization} if $T$ is lax
idempotent, and a \emph{colax localization} if $T$ is colax idempotent
(\cref{def:kz}).
\end{defi}

\begin{lem}
\label{lem:laxloc-criterion}
\gutentag{00CB}%
Let $f\colon x \to y$ and $u\colon y \to x$ be $1$-morphisms of an
$(\infty,2)$-category $\A$.
\begin{enumerate}
\item Given an invertible
$c\colon fu \simeq \mathrm{id}_{y}$,
we have $f \dashv u$ with counit $c$ if and only if there is
$\nu\colon \mathrm{id}_{x} \Rightarrow uf$
with $f\nu$ and $\nu u$ invertible.
\item Given an invertible
$\nu\colon \mathrm{id}_{x} \simeq uf$,
we have $f \dashv u$ with unit $\nu$ if and only if there is
$c\colon fu \Rightarrow \mathrm{id}_{y}$
with $cf$ and $uc$ invertible.
\end{enumerate}
\end{lem}

\begin{proof}
\gutentag{00CC}%
(1) The unit of an adjunction $f \dashv u$ with counit $c$ is such a $\nu$,
as $c$ is invertible, by
the triangle identities. Conversely, by \cite[Cor.~2.2]{bourke2025recognition},
a consequence of \cite[Lemma~2.1.11]{riehlverity2022elements}, a $\nu$ as in
(1) yields a $2$-morphism
$\mathrm{id}_{x} \Rightarrow uf$
satisfying the triangle identities with $c$.
\todo{2-categorical input: check that \cite[Cor.~2.2]{bourke2025recognition} works verbatim.} By \cref{thm:adj-lifting}(2), these extend to an adjunction
$f \dashv u$ with counit $c$. (2) is (1) in
$\A^{\mathrm{co}}$.
\end{proof}

\begin{lem}
\label{lem:laxloc-equiv}
\gutentag{00CD}%
The following are equivalent:
\begin{enumerate}
\item $L \dashv R$ is a lax localization;
\item
$L\eta \dashv \varepsilon L$
in $\func(\D,\C)$, with unit $\varphi^{-1}$;
\item
$R\varepsilon \dashv \eta R$
in $\func(\C,\D)$, with counit $\theta$;
\item
$L\eta_{d} \dashv \varepsilon_{Ld}$
with unit $\varphi_{d}^{-1}$, for
every $d \in \D$;
\item
$R\varepsilon_{c} \dashv \eta_{Rc}$
with counit $\theta_{c}$, for every
$c \in \C$.
\end{enumerate}
In particular, for a lax localization, $L\eta$ is coreflective in
$\func(\D,\C)$ and $\eta R$ is reflective in $\func(\C,\D)$
(\cref{def:reflection}), and so are their components $L\eta_{d}$ and
$\eta_{Rc}$.
\end{lem}

\begin{proof}
\gutentag{00CE}%
Evaluation preserves adjunctions, so (2) implies (4) and (3) implies (5).
Postcomposition with $R$ takes (2) to
$T\eta \dashv \mu$
with unit $(R\varphi)^{-1}$, and precomposition with $L$ takes (3) to
$\mu \dashv \eta T$
with counit $\theta L$; so (2) and (3) each imply (1), by \cref{def:kz} and
\cref{lem:kz-basic}(1).

Assume (1). Let $\upsilon$ be the unit of
$\mu \dashv \eta T$
(\cref{lem:kz-basic}(1)) and
$\delta\colon T\eta \Rightarrow \eta T$
the $2$-morphism $\upsilon T\eta$ followed by the unit law, as in
\cref{prop:kz-action}. By the triangle identities, as the counit of
$\mu \dashv \eta T$
is invertible, so are $\mu\upsilon$ and $\upsilon\eta T$. Hence $\mu\delta$ is
invertible, and so is $\delta_{d}\eta_{d}$: up to invertible $2$-morphisms, it
is $\upsilon_{d}$ whiskered with
$T\eta_{d} \circ \eta_{d} \simeq \eta_{Td} \circ \eta_{d}$.

(1)$\Rightarrow$(5). Let $x := Rc$ and
$a := R\varepsilon_{c}$,
so that
$\theta_{c}\colon a\eta_{x} \simeq \mathrm{id}_{x}$,
and let $\nu$ be the composite
\begin{equation*}
\mathrm{id}_{Tx} \simeq Ta \circ T\eta_{x} \xRightarrow{\ Ta\,\delta_{x}\ } Ta
\circ \eta_{Tx} \simeq \eta_{x} \circ a \ .
\end{equation*}
Then $\nu\eta_{x}$ is invertible as $\delta_{x}\eta_{x}$ is, and $a\nu$ is
invertible as $a\mu_{x}\delta_{x}$ is, since
$a \circ Ta \simeq a \circ \mu_{x}$
by $R$ applied to the naturality of $\varepsilon$ at $\varepsilon_{c}$.
So $a \dashv \eta_{x}$ with counit $\theta_{c}$ by
\cref{lem:laxloc-criterion}(1).

(1)$\Rightarrow$(4). The equivalence \eqref{eq:laxloc-adjunction}
$\Hom_{\C}(LTd,LTd) \simeq \Hom_{\D}(Td,TTd)$
sends
$L\eta_{d} \circ \varepsilon_{Ld}$
to
$T\eta_{d} \circ \mu_{d} \circ \eta_{Td} \simeq T\eta_{d}$
and $\mathrm{id}_{LTd}$ to $\eta_{Td}$, so $\delta_{d}$ corresponds to a
$2$-morphism
$\kappa\colon L\eta_{d} \circ \varepsilon_{Ld} \Rightarrow \mathrm{id}_{LTd}$.
By naturality of \eqref{eq:laxloc-adjunction},
$\varepsilon_{Ld}\kappa$
corresponds to $\mu_{d}\delta_{d}$ and $\kappa L\eta_{d}$ to
$\delta_{d}\eta_{d}$, so both are invertible, and
$L\eta_{d} \dashv \varepsilon_{Ld}$
with unit $\varphi_{d}^{-1}$ by \cref{lem:laxloc-criterion}(2).

(5)$\Rightarrow$(3). By
\cite[Cor.~5.2.12]{abellangagnahaugseng2024straightening},
$R\varepsilon$ is a left adjoint in $\func(\C,\D)$ if its components are and
its naturality squares are right Beck--Chevalley (\cref{lem:mates}(1)).
At $g\colon c \to c'$ the right mate is
\begin{equation*}
TRg \circ \eta_{Rc} \Longrightarrow \eta_{Rc'}R\varepsilon_{c'} \circ TRg \circ
\eta_{Rc} \simeq \eta_{Rc'} \circ Rg \circ R\varepsilon_{c} \circ \eta_{Rc}
\overset{\theta_{c}}{\simeq} \eta_{Rc'} \circ Rg \ ,
\end{equation*}
whose first $2$-morphism is the unit of
$R\varepsilon_{c'} \dashv \eta_{Rc'}$
at
$TRg \circ \eta_{Rc} \simeq \eta_{Rc'} \circ Rg$,
invertible by a triangle identity as the counit is invertible. So
$R\varepsilon \dashv G$
for some $G$. The $2$-morphism
$\eta R \Rightarrow G$
adjunct to $\theta$ has as components the comparisons between the right adjoints
$\eta_{Rc}$ and $G_{c}$ of $R\varepsilon_{c}$, so it is invertible, as
$2$-morphisms of $\func(\C,\D)$ are detected componentwise (see the proof of
\cref{lem:mates-refl}). Hence
$R\varepsilon \dashv \eta R$
with counit $\theta$.

(4)$\Rightarrow$(2). Likewise, at $g\colon d \to d'$ the right mate of the
naturality square of $L\eta$ is
\begin{equation*}
Lg \circ \varepsilon_{Ld} \xRightarrow{\ \varphi_{d'}^{-1}\ }
\varepsilon_{Ld'}L\eta_{d'} \circ Lg \circ \varepsilon_{Ld} \simeq
\varepsilon_{Ld'} \circ LTg \circ L\eta_{d} \circ \varepsilon_{Ld}
\Longrightarrow \varepsilon_{Ld'} \circ LTg \ ,
\end{equation*}
whose last $2$-morphism is the counit $\kappa_{d}$ of
$L\eta_{d} \dashv \varepsilon_{Ld}$
whiskered with
$\varepsilon_{Ld'} \circ LTg \simeq Lg \circ \varepsilon_{Ld}$;
it is invertible as
$\varepsilon_{Ld}\kappa_{d}$
is, by a triangle identity.
So $L\eta \dashv G'$ for some $G'$, and the $2$-morphism
$G' \Rightarrow \varepsilon L$
adjunct to $\varphi^{-1}$ is invertible, as its components are the comparisons
between the right adjoints $G'_{d}$ and $\varepsilon_{Ld}$ of $L\eta_{d}$.
\end{proof}

\begin{rmk}
\label{rmk:laxloc-literature}
\gutentag{00CF}%
For $2$-categories, Bourke \cite[\S 2.1]{bourke2025recognition} calls a
biadjunction lax-idempotent if (3) holds, and the equivalence of (1) and (3) is
\cite[Prop.~2.5.12]{stepan2025thesis}. For $(\infty,2)$-categories we know of no
reference (\cref{rmk:kz-literature}).
\end{rmk}

\gutentag{00D0}%
As $(-)^{\mathrm{co}}$ is an automorphism of $\twoCat$ with
$\iota_{1}\func(\A^{\mathrm{co}},\B^{\mathrm{co}}) \simeq
\iota_{1}\func(\A,\B)$
compatibly with composition (\cref{subsec:prelim-monads}), an adjunction
$L \dashv R$ is also an adjunction
$L^{\mathrm{co}} \dashv R^{\mathrm{co}}$
between $\D^{\mathrm{co}}$ and $\C^{\mathrm{co}}$, with the same unit, counit
and monad. By \cref{def:kz}, $L \dashv R$ is a colax localization if
and only if
$L^{\mathrm{co}} \dashv R^{\mathrm{co}}$
is a lax localization. We treat lax localizations, and will spell out the colax
case in \cref{rmk:laxloc-colax}.

\subsection{Local objects}
\label{subsec:laxloc-local}

Let
$P \subset \func(\Delta^{1},\D)$
be a locally full sub-$(\infty,2)$-category. We write the objects of $P$ as
$p\colon x \to y$, and its $1$-morphisms, which are commutative squares, as
$\tau = (a,b)\colon p \to p'$,
drawn with $p$ and $p'$ vertical as in \cref{lem:mates}(1).

\begin{defi}
\label{def:laxloc-local}
An object $z \in \D$ is \emph{$P$-local} if
\begin{enumerate}
\item[(O1)] for every object $p\colon x \to y$ of $P$, the functor
$p^{\natural}\colon \Hom_{\D}(y,z) \to \Hom_{\D}(x,z)$
has a fully faithful left adjoint $(p^{\natural})^{L}$, and
\item[(O2)] for every $1$-morphism
$\tau = (a,b)\colon (p\colon x \to y) \to (p'\colon x' \to y')$
of $P$, the commutative square
\begin{equation*}
\begin{tikzcd}[column sep=large]
\Hom_{\D}(y',z) \ar[r, "p'^{\natural}"] \ar[d, "b^{\natural}"'] &
\Hom_{\D}(x',z) \ar[d, "a^{\natural}"] \ar[dl, Rightarrow, shorten=4mm] \\
\Hom_{\D}(y,z) \ar[r, "p^{\natural}"'] & \Hom_{\D}(x,z)
\end{tikzcd}
\end{equation*}
filled by $\tau$ is left Beck--Chevalley, i.e.\
$(p^{\natural})^{L}a^{\natural} \Rightarrow b^{\natural}(p'^{\natural})^{L}$
is invertible.
\end{enumerate}
A $1$-morphism $g\colon z \to z'$ between $P$-local objects is
\emph{$P$-local} if
\begin{enumerate}
\item[(M)] for every object $p\colon x \to y$ of $P$, the square
\begin{equation*}
\begin{tikzcd}[column sep=large]
\Hom_{\D}(y,z) \ar[r, "p^{\natural}"] \ar[d, "g_{\natural}"'] &
\Hom_{\D}(x,z) \ar[d, "g_{\natural}"] \ar[dl, Rightarrow, shorten=4mm, "="'] \\
\Hom_{\D}(y,z') \ar[r, "p^{\natural}"'] & \Hom_{\D}(x,z')
\end{tikzcd}
\end{equation*}
given by associativity is left Beck--Chevalley, i.e.\
$(p^{\natural})^{L}g_{\natural} \Rightarrow g_{\natural}(p^{\natural})^{L}$
is invertible.
\end{enumerate}
Let
$\Pc_{\emptyset,P} \subset \D$
be the locally full sub-$(\infty,2)$-category of $P$-local objects and
$1$-morphisms.
\end{defi}

By \cref{lem:mate-correspondence}(3), $P$-local $1$-morphisms are closed under
composition and contain the identities, and both classes are closed under
equivalence, so $\Pc_{\emptyset,P}$ exists (\cref{subsec:prelim-conventions}).
We will see in \cref{thm:local-pullback} that it is the case $I = Q =
\emptyset$ of \cref{def:IP-local}, which explains the notation.
If the $1$-morphisms of
$P$ are just the equivalences, the $P$-local objects and $1$-morphisms are the
left Kan injective ones of \cite[Def.~1.2,
Rem.~1.3]{dilibertilobbiasousa2022kz}. We
write $\{\sigma\}$ for the locally full sub-$(\infty,2)$-category of
$\func(\Delta^{1},\D)$ generated by a $1$-morphism or a commutative square
$\sigma$, and $\{s_{j}\}_{j}$, or just $\{s_{j}\}$,
for the one whose objects are the $1$-morphisms
$s_{j}$ and whose $1$-morphisms are the equivalences. Then $\Pc_{\emptyset,P}$
is the intersection of the
$\Pc_{\emptyset,\{p\}}$
and
$\Pc_{\emptyset,\{\tau\}}$
over the objects $p$ and $1$-morphisms $\tau$ of $P$.

As $(p^{\natural})^{L}$ is fully faithful, the counit of
$(p^{\natural})^{L} \dashv p^{\natural}$
is invertible exactly at the objects of its essential image. By
\cref{constr:mates}, the left mates of (O2) and (M) are, at $m$, this counit at
$b^{\natural}(p'^{\natural})^{L}m$,
resp.\ at
$g_{\natural}(p^{\natural})^{L}m$,
composed with invertible $2$-morphisms. So, granting (O1), (O2) says that
$b^{\natural}$ carries the essential image of $(p'^{\natural})^{L}$ into that of
$(p^{\natural})^{L}$, and (M) says that $g_{\natural}$ carries the essential
image of $(p^{\natural})^{L}$ on $\Hom_{\D}(y,z)$ into that on
$\Hom_{\D}(y,z')$.

\begin{lem}
\label{lem:laxloc-closure}
\leavevmode
\begin{enumerate}
\item A left adjoint between $P$-local objects is $P$-local.
\item If $d$ is $P$-local and $i\colon e \to d$ is reflective, then $e$ and
$i^{L}$ are $P$-local.
\item Let $i\colon e \to d$ be reflective and $f\colon e \to e'$ a $1$-morphism
such that $i^{L}$ and $fi^{L}$ are $P$-local. Then $f$ is $P$-local.
\end{enumerate}
\end{lem}

\begin{proof}
(1) Let
$g \dashv g^{R}$
with $g\colon z \to z'$. The square of (M) has the right mate
$p^{\natural}(g^{R})_{\natural} \Rightarrow (g^{R})_{\natural}p^{\natural}$,
which is invertible, being $p^{\natural}$ applied to a triangle identity. So the
square is left Beck--Chevalley by \cref{lem:mates-conjugate}.

(2) Let $p\colon x \to y$ be an object of $P$, $(p^{\natural})^{L}$ the left
adjoint of $p^{\natural}$ on $\Hom_{\D}(-,d)$, and
$\lambda := i^{L}_{\natural}(p^{\natural})^{L}i_{\natural}\colon
\Hom_{\D}(x,e) \to \Hom_{\D}(y,e)$.
As $i_{\natural}$ is fully faithful and commutes with $p^{\natural}$, there
are natural equivalences
\begin{equation*}
\mapp(\lambda m,n) \simeq \mapp\bigl((p^{\natural})^{L}im,in\bigr) \simeq
\mapp(im,inp) \simeq \mapp(m,np) \ ,
\end{equation*}
so
$\lambda \dashv p^{\natural}$
\cite[Prop.~5.2.2.8]{htt}. Under $i_{\natural}$,
its unit at $m$ is the unit of $(p^{\natural})^{L}$ at $im$ followed by the unit
of
$i^{L} \dashv i$
at
$(p^{\natural})^{L}(im) \circ p \simeq im$;
both are invertible, the latter as $im$ lies in the essential image of
$i_{\natural}$. So (O1) holds for $e$. For $\tau = (a,b)$ in $P$, (O2) for $d$
gives
\begin{equation*}
b^{\natural}\lambda' =
i^{L}_{\natural}b^{\natural}(p'^{\natural})^{L}i_{\natural}
\simeq i^{L}_{\natural}(p^{\natural})^{L}a^{\natural}i_{\natural} = \lambda
a^{\natural} \ ,
\end{equation*}
with $\lambda'$ defined like $\lambda$ for $p'$. So $b^{\natural}$ carries the
essential image of $\lambda'$ into that of $\lambda$, which is (O2) for $e$.
Now $i^{L}$ is $P$-local by (1).

(3) The objects $d$, $e$ and $e'$ are $P$-local, being sources and targets of
$P$-local $1$-morphisms. By \cref{lem:mate-correspondence}(3),
the left mate of (M) for $fi^{L}$ is the
pasting of those for $i^{L}$ and $f$. As the former two are invertible, the
left mate for $f$ is invertible after precomposition with $i^{L}_{\natural}$,
which is essentially surjective as
$i^{L}_{\natural}i_{\natural} \simeq \mathrm{id}$.
So $f$ is $P$-local.
\end{proof}

\begin{lem}
\label{lem:laxloc-repr}
A $1$-morphism $s$ of $\D$ is coreflective if and only if
$\Pc_{\emptyset,\{s\}} = \D$.
A commutative square $\sigma$ between coreflective $1$-morphisms, drawn with
these vertical, is right Beck--Chevalley if and only if
$\Pc_{\emptyset,\{\sigma\}} = \D$.
\end{lem}

\begin{proof}
The enriched Yoneda embedding
$h\colon \D\opcat \to \func(\D,\Cat)$,
$d \mapsto \Hom_{\D}(d,-)$,
is fully faithful \cite{hinich2020yoneda,heine2023enriched}. It takes an
adjunction $l \dashv r$ of $\D$ to
$r^{\natural} \dashv l^{\natural}$
with the same unit and counit, and every adjunction between representables
arises this way, as adjunctions are detected in the homotopy $2$-category
(\cref{subsec:prelim-adjunctions}). So $s\colon x \to y$ is coreflective if and
only if
$s^{\natural}\colon h(y) \to h(x)$
is a right adjoint with invertible unit. By
\cite[Cor.~5.2.12]{abellangagnahaugseng2024straightening}, $s^{\natural}$ is a
right adjoint if and only if its components are and its transposed naturality
squares, the squares of (M), are left Beck--Chevalley; and a $2$-morphism of
$\func(\D,\Cat)$ is invertible if and only if its components are. This proves
the first claim. For
$\sigma = (a,b)\colon s \to s'$,
$h$ takes the right mate
$as^{R} \Rightarrow s'^{R}b$
to the left mate
$(s^{\natural})^{L}a^{\natural} \Rightarrow b^{\natural}(s'^{\natural})^{L}$
of $h(\sigma)$, whose components are the $2$-morphisms of (O2), and $h$
reflects invertibility.
\end{proof}

\begin{prop}
\label{prop:laxloc-adjunction}
Let $s$ be a $1$-morphism of $\D$. Then $Ls$ is coreflective if and only if
every object $Rc$ and every $1$-morphism $Rg$ is $\{s\}$-local. If $\sigma$ is
a commutative square between such $1$-morphisms, then $L\sigma$ is right
Beck--Chevalley if and only if every $Rc$ is $\{\sigma\}$-local.
\end{prop}

\begin{proof}
By naturality of \eqref{eq:laxloc-adjunction}, the squares of
\cref{def:laxloc-local} for $Ls$, $L\sigma$ and an object $c$ or a $1$-morphism
$g$ of $\C$ are equivalent to those for $s$, $\sigma$ and $Rc$, $Rg$. So this
is \cref{lem:laxloc-repr} in $\C$.
\end{proof}

\subsection{The modules of a lax localization}
\label{subsec:laxloc-description}

\gutentag{00DB}%
Let $L \dashv R$ be a lax localization. By
\cref{thm:kz-modules,thm:kz-module-maps}, the $T$-modules form the locally full
sub-$(\infty,2)$-category
$\D_{T} \subset \D$
whose objects are the $x$ with $\eta_{x}$ reflective and whose $1$-morphisms
are those with left Beck--Chevalley unit square, and
$\iota_{1}\D_{T} \simeq \modcat{T}{\D}$.
For $x \in \D_{T}$ we write
$a_{x}\colon Tx \to x$
for the left adjoint of $\eta_{x}$, the action.
Let
$\Sigma_{\eta} \subset \func(\Delta^{1},\D)$
be the locally full sub-$(\infty,2)$-category generated by the $\eta_{d}$ and
their naturality squares
$(g,Tg)\colon \eta_{d} \to \eta_{d'}$
for $g\colon d \to d'$.

\begin{defi}
\label{def:closed-loc}
\gutentag{00DC}%
A locally full sub-$(\infty,2)$-category
$\B \subset \D$
is \emph{closed under localizations} if:
\begin{enumerate}
\item[(R1)] if $d$ is in $\B$ and $i\colon e \to d$ is reflective, then $e$ and
$i^{L}$ are in $\B$;
\item[(R2)] if $i\colon e \to d$ is reflective and $i^{L}$ and $fi^{L}$ are in
$\B$, then $f$ is in $\B$.
\end{enumerate}
\end{defi}

\gutentag{00DD}%
By \cref{lem:laxloc-closure}, every $\Pc_{\emptyset,P}$ is closed under
localizations.

\begin{thm}
\label{thm:laxloc-local}
\gutentag{00DE}%
Let $L \dashv R$ be a lax localization.
\begin{enumerate}
\item $R$ factors through $\D_{T}$.
\item
$\D_{T} = \Pc_{\emptyset,\Sigma_{\eta}} =
\Pc_{\emptyset,\{\eta_{d}\}_{d \in \D}}$.
\item $\D_{T}$ is the smallest locally full sub-$(\infty,2)$-category of $\D$
that contains the free modules $Td$ and the $1$-morphisms $Tg$ and is closed
under localizations. In particular, an object of $\D$ lies in $\D_{T}$ if and
only if
it admits
a reflective $1$-morphism to some $Rc$.
\end{enumerate}
\end{thm}

\begin{proof}
\gutentag{00DF}%
(1) By \cref{lem:laxloc-equiv}(3),
$R\varepsilon \dashv \eta R$
is a reflection in $\func(\C,\D)$. By \cref{lem:mates-refl}(2), every
$\eta_{Rc}$ is reflective and the naturality squares of $\eta R$, which are the
unit squares of the $Rg$, are left Beck--Chevalley. So $R$ factors through
$\D_{T}$ \cite[Cor.~2.6.6]{abellangagnahaugseng2024straightening}.

(2) By \cref{lem:laxloc-equiv}(2) and \cref{lem:mates-refl}(2), every
$L\eta_{d}$ is coreflective and $L$ takes the squares $(g,Tg)$ to right
Beck--Chevalley squares, so every $Rc$ and $Rg$ is $\Sigma_{\eta}$-local by
\cref{prop:laxloc-adjunction}. We show
$\D_{T} \subset \Pc_{\emptyset,\Sigma_{\eta}} \subset
\Pc_{\emptyset,\{\eta_{d}\}} \subset \D_{T}$,
where the middle inclusion is clear. For $x \in \D_{T}$,
\cref{lem:laxloc-closure}(2)
for the reflective
$\eta_{x}\colon x \to Tx = RLx$
shows that $x$ and $a_{x}$ are $\Sigma_{\eta}$-local. For
$f\colon x \to y$
in $\D_{T}$, the left mate of its unit square is an equivalence
$a_{y} \circ Tf \simeq f \circ a_{x}$,
so $fa_{x}$ is $\Sigma_{\eta}$-local, and so is $f$ by
\cref{lem:laxloc-closure}(3).

For the last inclusion, let $\Lambda^{z}_{d}$ denote the left adjoint of
$\eta_{d}^{\natural}\colon \Hom_{\D}(Td,z) \to \Hom_{\D}(d,z)$
for $\{\eta_{d}\}$-local $z$. For $z = Ty$, \eqref{eq:laxloc-adjunction}
identifies $\eta_{d}^{\natural}$ with
$(L\eta_{d})^{\natural}\colon \Hom_{\C}(LTd,Ly) \to \Hom_{\C}(Ld,Ly)$,
whose left adjoint is
$\varepsilon_{Ld}^{\natural}$
by
\cref{lem:laxloc-equiv}(4). As
$m \circ \varepsilon_{Ld}$
corresponds to
$Rm \circ \mu_{d} \circ \eta_{Td} \simeq Rm$,
the essential image of $\Lambda^{Ty}_{d}$ consists of the $Rm$ for
$m\colon Ld \to Ly$.

Let $x$ be $\{\eta_{d}\}$-local, let
$a := \Lambda^{x}_{x}(\mathrm{id}_{x})$,
and let
$c\colon a\eta_{x} \simeq \mathrm{id}_{x}$
be the inverse of the unit. As
$\mathrm{id}_{Tx} = R(\mathrm{id}_{Lx})$,
we have
$\mathrm{id}_{Tx} \simeq \Lambda^{Tx}_{x}(\eta_{x})$,
so the $2$-morphism
$\eta_{x}c^{-1}\colon \eta_{x} \Rightarrow \eta_{x}^{\natural}(\eta_{x}a)$
has an adjunct
$\nu\colon \mathrm{id}_{Tx} \Rightarrow \eta_{x}a$.
Then $\nu\eta_{x}$ is invertible, as the unit of
$\Lambda^{Tx}_{x} \dashv \eta_{x}^{\natural}$
is. Both $a$ and
$a\eta_{x}a \simeq a$
lie in the essential image of $\Lambda^{x}_{x}$, on which $\eta_{x}^{\natural}$
is fully faithful, and
$\eta_{x}^{\natural}(a\nu) = a\nu\eta_{x}$
is invertible; so $a\nu$ is invertible, and
$a \dashv \eta_{x}$
with counit $c$ by \cref{lem:laxloc-criterion}(1). So $x \in \D_{T}$ and
$a_{x} = a$.

Let $f\colon x \to y$ be $\{\eta_{d}\}$-local. The left adjoints $a_{x}$ and
$a_{y}$ are $\{\eta_{d}\}$-local by \cref{lem:laxloc-closure}(1). As
$a_{x} = \Lambda^{x}_{x}(\mathrm{id}_{x})$
and $Tf = R(Lf)$ lies in the essential image of $\Lambda^{Ty}_{x}$,
the criterion
after \cref{def:laxloc-local} shows that $f \circ a_{x}$ and
$a_{y} \circ Tf$
lie in the essential image of $\Lambda^{y}_{x}$. So the left mate
$a_{y} \circ Tf \Rightarrow f \circ a_{x}$
of the unit square of $f$ is invertible if its whiskering with $\eta_{x}$ is.
By \cref{constr:mates}, this whiskering is
$a_{y}Tf\,u_{x}\eta_{x}$
composed
with invertible $2$-morphisms, where $u_{x}$ is the unit of
$a_{x} \dashv \eta_{x}$,
and $u_{x}\eta_{x}$ is invertible by a triangle identity. So $f$ is in
$\D_{T}$.

(3) By (1) and (2), $\D_{T}$ contains the $Rc$ and $Rg$, in particular the $Td$
and $Tg$, and is closed under localizations. Conversely, let $\B$ be such. For
$x \in \D_{T}$, (R1) for
$\eta_{x}\colon x \to Tx$
shows that $x$ and $a_{x}$ are in $\B$, and for $f\colon x \to y$ in $\D_{T}$,
(R2) for
$fa_{x} \simeq a_{y} \circ Tf$,
which is in $\B$, shows that $f$ is in $\B$. The last claim follows from (R1)
for $\D_{T}$, as
$\eta_{x}\colon x \to RLx$
is reflective for $x \in \D_{T}$.
\end{proof}

\begin{rmk}
\label{rmk:laxloc-local}
\gutentag{00E0}%
For $2$-categories, the equality
$\D_{T} = \Pc_{\emptyset,\{\eta_{d}\}}$
is \cite[Thm.~2.7]{dilibertilobbiasousa2022kz}, based on Marmolejo--Wood
\cite{marmolejowood2012kan}, and \cref{lem:laxloc-closure}(2), (3) are
contained in \cite[Lemma~2.6]{dilibertilobbiasousa2022kz}. By (3), $\D_{T}$ is
the closure of the free modules under localizations, just as the local objects
of a Bousfield localization are the closure of the free modules under
equivalences.
\end{rmk}

\begin{prop}
\label{prop:laxloc-normalize}
\gutentag{00E1}%
Let $L \dashv R$ be a lax localization, and
$\iota\colon \D_{T} \to \D$
the inclusion. Then
$T \simeq \iota L_{T}$
for a functor
$L_{T}\colon \D \to \D_{T}$,
and
$L_{T} \dashv \iota$
with unit $\eta$, and its counit has components $a_{x}$. This adjunction is a
lax localization, whose monad has underlying functor $T$ and unit $\eta$.
\end{prop}

\begin{proof}
\gutentag{00E2}%
By \cref{thm:laxloc-local}(1), $T = RL$ factors as $\iota L_{T}$. The
components $\eta_{x}$ of
$\eta\iota\colon \iota \Rightarrow T\iota$
are reflective and its naturality squares, the unit squares of the
$1$-morphisms of $\D_{T}$, are left Beck--Chevalley. So by
\cite[Cor.~5.2.12]{abellangagnahaugseng2024straightening} (\cref{lem:mates}(2))
it has a left adjoint $a$ in $\func(\D_{T},\D)$, with components $a_{x}$ and
invertible counit. The $a_{x}$ are in $\D_{T}$ by \cref{lem:laxloc-closure}(1)
and \cref{thm:laxloc-local}(2), so $a$ lifts to
$L_{T}\iota \Rightarrow \mathrm{id}_{\D_{T}}$,
as $\func(\D_{T},-)$ preserves locally full inclusions
(\cref{subsec:prelim-conventions}). For the triangle identities,
$a \circ \eta\iota \simeq \mathrm{id}_{\iota}$
is the counit of
$a \dashv \eta\iota$,
and
$\iota(aL_{T} \circ L_{T}\eta) \simeq aL_{T} \circ T\eta \simeq \mu \circ T\eta
\simeq \mathrm{id}_{T}$,
as $aL_{T}$ and $\mu$ are both left adjoint to $\eta T$
(\cref{lem:kz-basic}(1)); this lifts along the locally fully faithful
$\iota_{\natural}\colon \func(\D,\D_{T}) \to \func(\D,\D)$.
So $\eta$ and the lift of $a$ are the unit and counit of an adjunction
$L_{T} \dashv \iota$
(\cref{subsec:prelim-adjunctions}).
Finally, \cref{lem:laxloc-equiv}(5) holds for
$L_{T} \dashv \iota$,
as
$a_{x} \dashv \eta_{x}$
with counit the triangle identity.
\end{proof}

\begin{defi}
\label{def:laxloc-normalized}
\gutentag{00E3}%
A lax localization is \emph{normalized} if $R$ factors through an equivalence
$\C \simeq \D_{T}$.
\end{defi}

\gutentag{00E4}%
By \cref{thm:laxloc-local}, a lax localization is normalized if and only if $R$
is a locally full inclusion whose image is closed under localizations, and
\cref{prop:laxloc-normalize} replaces every lax localization by a normalized one
with the same monad functor and unit. Normalization is not automatic:

\begin{exmp}
\label{ex:laxloc-initial}
\gutentag{00E5}%
If $\C$ has an object $0$ with
$\Hom_{\C}(0,c) \simeq \mathrm{pt}$
for all $c$, then
$0\colon \mathrm{pt} \to \C$
is left adjoint to
$\C \to \mathrm{pt}$,
with $T = \mathrm{id}$. This is a lax localization with
$\D_{T} = \mathrm{pt}$,
normalized only if
$\C \simeq \mathrm{pt}$.
\end{exmp}

\gutentag{00E6}%
Let
$\Sat_{L} \subset \func(\Delta^{1},\D)$
be the preimage under $L$ of the image of
$\ell^{*}\colon \func(\Corefl,\C) \to \func(\Delta^{1},\C)$:
its objects are the $1$-morphisms $f$ with $Lf$ coreflective, and its
$1$-morphisms are the commutative squares $\sigma$ between them,
drawn with these
vertical, with $L\sigma$ right Beck--Chevalley (\cref{lem:mates-refl}). It is
locally full, as locally full inclusions are stable under pullback.

\begin{prop}
\label{prop:laxloc-class}
\gutentag{00E7}%
Let $L \dashv R$ be a lax localization.
\begin{enumerate}
\item
$\Sigma_{\eta} \subset \Sat_{L}$.
\item A $1$-morphism $f$ is in $\Sat_{L}$ if and only if
$\D_{T} \subset \Pc_{\emptyset,\{f\}}$,
and a commutative square $\sigma$ between $1$-morphisms in $\Sat_{L}$ is an
$\Sat_{L}$-square if and only if
$\D_{T} \subset \Pc_{\emptyset,\{\sigma\}}$.
\item
$\Pc_{\emptyset,\Sat_{L}} = \D_{T}$,
and $\Sat_{L}$ is the largest locally full sub-$(\infty,2)$-category $P$ of
$\func(\Delta^{1},\D)$ with
$\D_{T} \subset \Pc_{\emptyset,P}$.
\end{enumerate}
\end{prop}

\begin{proof}
\gutentag{00E8}%
(1) was shown in the proof of \cref{thm:laxloc-local}(2). (2) By
\cref{prop:laxloc-adjunction}, $f$ is in $\Sat_{L}$ if and only if every $Rc$
and $Rg$ lies in
$\Pc_{\emptyset,\{f\}}$,
and as
$\Pc_{\emptyset,\{f\}}$
is closed under localizations, this holds if and only if
$\D_{T} \subset \Pc_{\emptyset,\{f\}}$
by \cref{thm:laxloc-local}(1), (3). Likewise for $\sigma$.
(3) By (2),
$\D_{T} \subset \Pc_{\emptyset,\Sat_{L}}$,
and by (1) and \cref{thm:laxloc-local}(2),
$\Pc_{\emptyset,\Sat_{L}} \subset \Pc_{\emptyset,\Sigma_{\eta}} = \D_{T}$.
If
$\D_{T} \subset \Pc_{\emptyset,P}$,
then
$\D_{T} \subset \Pc_{\emptyset,\{\sigma\}}$
for every object and $1$-morphism $\sigma$ of $P$, so
$P \subset \Sat_{L}$
by (2).
\end{proof}

\begin{rmk}
\label{rmk:laxloc-saturated}
\gutentag{00E9}%
By \cref{prop:laxloc-class}(2) and \cref{lem:saturation-explicit}, which we will
prove in \cref{subsec:2dloc-saturation},
$\Sat_{L} = \Sat^{\mathrm{corefl}}_{\D_{T}}$,
so $\Sat_{L}$ is coreflectively saturated by \cref{thm:saturation-saturated},
without any presentability; and
$\Sat_{L} = \overline{\Sigma_{\eta}} = \overline{P}$
for every $P$ with
$\Pc_{\emptyset,P} = \D_{T}$
(\cref{def:saturation}). For an arbitrary adjunction,
\cref{prop:laxloc-adjunction}
shows in the same way that $\Sat_{L}$ is the coreflective saturation of the
locally full sub-$(\infty,2)$-category generated by the image of $R$, so it is
coreflectively saturated as well.
\end{rmk}

\subsection{Recognizing lax localizations}
\label{subsec:laxloc-recognition}

For a $1$-morphism $f\colon d \to e$ of $\D$ we call the lax pullback
$(e \downarrow f) := (\mathrm{id}_{e} \downarrow f)$
(\cref{subsec:prelim-adjunctions}), if it exists, the \emph{colax limit} of $f$,
following \cite[\S 2.3]{bourke2025recognition}. We write
$\bar{p}\colon (e \downarrow f) \to e$,
$\bar{q}\colon (e \downarrow f) \to d$
and
$\bar{\alpha}\colon \bar{p} \Rightarrow f\bar{q}$
for its structure, so that
$h \mapsto (\bar{p}h,\bar{q}h,\bar{\alpha}h)$
is an equivalence
$\Hom_{\D}(t,(e \downarrow f)) \simeq (\Hom_{\D}(t,e) \downarrow f_{\natural})$,
naturally in $t$. For example, $\Cat$ has the colax limits of all
$1$-morphisms, given by comma $\infty$-categories.

\begin{lem}
\label{lem:laxloc-colax-limits}
Let
$P \subset \func(\Delta^{1},\D)$
be locally full, and let $f\colon d \to e$ be a $1$-morphism between $P$-local
objects that has a colax limit. Then $(e \downarrow f)$, $\bar{p}$ and $\bar{q}$
are $P$-local.
\end{lem}

\begin{proof}
Let $p\colon x \to y$ be an object of $P$. Under the equivalences above,
$p^{\natural}$ on
$\Hom_{\D}(-,(e \downarrow f))$
is the functor
$(\Hom_{\D}(y,e) \downarrow f_{\natural}) \to (\Hom_{\D}(x,e) \downarrow
f_{\natural})$,
$(m_{1},m_{2},\kappa) \mapsto (m_{1}p,m_{2}p,\kappa p)$.
Let $\lambda_{e}$ and $\lambda_{d}$ be the left adjoints of $p^{\natural}$ on
$\Hom_{\D}(-,e)$ and $\Hom_{\D}(-,d)$, and
$\zeta\colon \lambda_{e}f_{\natural} \Rightarrow f_{\natural}\lambda_{d}$
the left mate of (M) for $f$, which need not be invertible. Then
\begin{equation*}
\Lambda(m_{1},m_{2},\kappa) := \bigl(\lambda_{e}m_{1},\,\lambda_{d}m_{2},\,
\zeta_{m_{2}} \circ \lambda_{e}\kappa\bigr)
\end{equation*}
is left adjoint to $p^{\natural}$, with unit given componentwise by the units of
$\lambda_{e}$ and $\lambda_{d}$: the mapping spaces of a comma $\infty$-category
are pullbacks of mapping spaces of the factors \todo{Reference.}, the adjunction
equivalences of $\lambda_{e}$ and $\lambda_{d}$ are compatible with these
pullbacks, as $\zeta$ corresponds under
$\lambda_{e} \dashv p^{\natural}$
to $f_{\natural}$ applied to the unit of $\lambda_{d}$ (\cref{constr:mates} and
a triangle identity), and \cite[Prop.~5.2.2.8]{htt} applies. As the unit is
invertible, $\Lambda$ is fully faithful, which is (O1). The counit is
componentwise as well, so the essential image of $\Lambda$ consists of the
$(n_{1},n_{2},\kappa)$ with $n_{1}$ and $n_{2}$ in the essential images of
$\lambda_{e}$ and $\lambda_{d}$. As $b^{\natural}$, $\bar{p}_{\natural}$ and
$\bar{q}_{\natural}$ act componentwise, the criterion after
\cref{def:laxloc-local} gives (O2) for $(e \downarrow f)$ from (O2) for $d$ and
$e$, and (M) for $\bar{p}$ and $\bar{q}$.
\end{proof}

\begin{thm}
\label{thm:laxloc-bourke}
Let $L \dashv R$ be an adjunction with $R$ locally fully faithful. Suppose that
for every $c \in \C$ the unit $\eta_{Rc}$ has a colax limit of the form $RZ$,
with
$\bar{p} = R(p_{Z})$
and
$\bar{q} = R(q_{Z})$
for some
$p_{Z}\colon Z \to LRc$
and
$q_{Z}\colon Z \to c$.
Then $L \dashv R$ is a lax localization.
\end{thm}

\begin{proof}
We follow \cite[Thm.~3.1]{bourke2025recognition} and verify
\cref{lem:laxloc-equiv}(5). Let $c \in \C$, $x := Rc$ and
$a := R\varepsilon_{c}$.
The colax limit provides
$h\colon x \to RZ$
with
$\bar{p}h \simeq \eta_{x}$,
$\bar{q}h \simeq \mathrm{id}_{x}$
and $\bar{\alpha}h$ identified with
$\mathrm{id}_{\eta_{x}}$.
By \eqref{eq:laxloc-adjunction},
$h \simeq Rh' \circ \eta_{x}$
for some
$h'\colon Lx \to Z$,
and as
$R(p_{Z}h')\eta_{x} \simeq \eta_{x}$
and
$R(q_{Z}h')\eta_{x} \simeq \mathrm{id}_{x} \simeq R(\varepsilon_{c})\eta_{x}$,
\eqref{eq:laxloc-adjunction} also gives
$p_{Z}h' \simeq \mathrm{id}_{Lx}$
and
$q_{Z}h' \simeq \varepsilon_{c}$.
So
$\bar{p}Rh' \simeq \mathrm{id}_{Tx}$
and
$\bar{q}Rh' \simeq a$,
and $\bar{\alpha}Rh'$ becomes a $2$-morphism
$\nu\colon \mathrm{id}_{Tx} \Rightarrow \eta_{x}a$.
Then $\nu\eta_{x}$ is invertible, as $\bar{\alpha}h$ is. As $R$ is locally fully
faithful, the $2$-morphism
$\theta_{c}a \circ a\nu\colon a \Rightarrow a$
is $R\gamma$ for some
$\gamma\colon \varepsilon_{c} \Rightarrow \varepsilon_{c}$,
and
$R\gamma\,\eta_{x} = \theta_{c}a\eta_{x} \circ a\nu\eta_{x}$
is invertible. By \eqref{eq:laxloc-adjunction}, $\gamma$ is invertible, hence
so is $a\nu$, and
$a \dashv \eta_{x}$
with counit $\theta_{c}$ by \cref{lem:laxloc-criterion}(1).
\end{proof}

\begin{rmk}
\label{rmk:laxloc-bourke}
Bourke assumes the colax limits of all $1$-morphisms
$Rc \to Rc'$
in the form of \cref{thm:laxloc-bourke}. The proof only uses those of the
units, only that
$\theta_{c}a \circ a\nu$
lies in the image of $R$, and only the
one-dimensional universal property of the colax limits, i.e.\ that the functors
$\Hom_{\D}(t,(e \downarrow f)) \to (\Hom_{\D}(t,e) \downarrow f_{\natural})$
are essentially surjective \cite[Rem.~3.2]{bourke2025recognition}.
\end{rmk}

\begin{cor}
\label{cor:laxloc-P}
Let
$P \subset \func(\Delta^{1},\D)$
be locally full such that
$\iota\colon \Pc_{\emptyset,P} \to \D$
has a left adjoint $L$, and suppose that every $1$-morphism of $\D$ between
$P$-local objects has a colax limit. Then
$L \dashv \iota$
is a normalized lax localization. In particular
$\Pc_{\emptyset,P} = \D_{T} = \Pc_{\emptyset,\Sigma_{\eta}}$,
$\eta_{x}$ is reflective for every $P$-local $x$, and
$P \subset \Sat_{L}$.
\end{cor}

\begin{proof}
$\iota$ is a locally full inclusion. For $P$-local $c$,
$\eta_{c}\colon c \to Tc$
is a $1$-morphism between $P$-local objects, so its colax limit and the
projections lie in $\Pc_{\emptyset,P}$ by \cref{lem:laxloc-colax-limits}, and
\cref{thm:laxloc-bourke} applies. The lax localization is normalized by
\cref{lem:laxloc-closure} and the paragraph after
\cref{def:laxloc-normalized}; the rest is \cref{thm:laxloc-local}(2),
\cref{lem:laxloc-equiv}(5) and \cref{prop:laxloc-class}(3).
\end{proof}

\begin{cor}
\label{cor:laxloc-classes}
Suppose that $\D$ has colax limits of all $1$-morphisms. Then
$\C \mapsto \Sat_{L}$
and
$P \mapsto \Pc_{\emptyset,P}$
are mutually inverse bijections between
\begin{enumerate}
\item the locally full sub-$(\infty,2)$-categories $\C \subset \D$ whose
inclusion $\iota$ has a left adjoint $L$ such that
$L \dashv \iota$
is a normalized lax localization, and
\item the locally full sub-$(\infty,2)$-categories
$P \subset \func(\Delta^{1},\D)$
such that
$\iota\colon \Pc_{\emptyset,P} \to \D$
has a left adjoint $L$ and $P = \Sat_{L}$.
\end{enumerate}
\end{cor}

\begin{proof}
For $\C$ as in (1), $\C = \D_{T}$, so
$\Pc_{\emptyset,\Sat_{L}} = \C$
by \cref{prop:laxloc-class}(3), and $\Sat_{L}$ is as in (2). For $P$ as in (2),
$L \dashv \iota$
is a normalized lax localization by \cref{cor:laxloc-P}, so
$\Pc_{\emptyset,P}$
is as in (1), with class $\Sat_{L} = P$.
\end{proof}

By \cref{rmk:laxloc-saturated} and \cref{cor:laxloc-P}, the condition
$P = \Sat_{L}$
in (2) says that
$P = \overline{P}$
for the saturation that will be defined in \cref{def:saturation}. In one
dimension, the analogous bijection is between reflective full subcategories and
the classes $W$ of morphisms whose local objects form a reflective full
subcategory and that equal the class of $W$-local equivalences, and
for a presentable $\infty$-category and a set $S$, the $S$-local equivalences
form the strongly saturated class generated by $S$ \cite[5.5.4.15]{htt}. We do
not know whether $\overline{P}$
is generated by $P$ under the operations of \cref{def:saturated}
(\cref{rmk:sat-open}(2)).

\begin{rmk}
\label{rmk:laxloc-presentable}
Let $\D \in \tprl$ and let $P$ be essentially small. Then
$\iota\colon \Pc_{\emptyset,P} \to \D$
will have a left adjoint by \cref{thm:2dloc-local}, and $\D$ has the colax
limits
$(e \downarrow f) \simeq e^{\Delta^{1}} \times_{e} d$
of all $1$-morphisms $f\colon d \to e$: limits in $\iota_{1}\D$ are
$(\infty,2)$-categorical by the argument that we will give for colimits in
\cref{subsec:2dloc-presentable}, as
$\mapp_{\D}(K \otimes t,-)$
preserves limits. So \cref{cor:laxloc-P} applies:
$L \dashv \iota$
is a normalized lax localization, and
$\Sigma_{\eta} \subset \overline{P}$,
as
$\Pc_{\emptyset,P} = \Pc_{\emptyset,\Sigma_{\eta}}$.
The same holds under the hypotheses of \cref{cor:laxloc-P}. There, this answers
the question of \cref{rmk:sat-open}(1), gives the hypothesis of
\cref{thm:sat-units}(3) (see \cref{rmk:sat-units-kz}), and shows that $T$ is
lax idempotent in \cref{conj:kleisli}, as asked in \cref{rmk:kleisli}.
\todo{Update these remarks when \cref{sec:2dloc} is restructured.}
\end{rmk}

\begin{rmk}
\label{rmk:laxloc-adjoining}
By \cref{lem:laxloc-equiv}(2) and \cref{prop:adjoin-univ},
which we will prove in
\cref{subsec:2dloc-adjoining}, a lax localization $L$ factors as
$\D \xrightarrow{\gamma} \D\corefl{\Sigma_{\eta}} \xrightarrow{\Phi} \C$,
and $\Phi$ takes values in the full sub-$(\infty,2)$-category
$K \subset \C$
on the objects $Ld$, as $\gamma$ is essentially surjective. The statement of
\cref{conj:kleisli}, that $\Phi$ is fully faithful, makes sense for every lax
localization, and the evidence in \cref{rmk:kleisli} applies verbatim. Granting
it, $L$ is a formal adjoining of coreflections, i.e.\ $\Phi$ is an equivalence,
if and only if every object of $\C$ is equivalent to some $Ld$. For normalized
lax localizations this fails in general: by \cref{thm:laxloc-local}(3),
$\D_{T}$ is the closure of $K$ under (R1) and (R2), and in
\cref{ex:local-vs-pushout}(2) (with $\D = \Cat$ and
$P = \{\emptyset \to \mathrm{pt}\}$)
this closure is strictly larger than $K$. In the presentable setting,
\cref{thm:2dloc-local} identifies the presentable formal adjoining
$\D\corefl{P}$ with $\Pc_{\emptyset,P}$, a normalized lax localization by
\cref{rmk:laxloc-presentable}; so the normalized lax localizations of this form
are those with
$\D_{T} = \Pc_{\emptyset,P}$
for some essentially small $P$. We do not know which normalized lax
localizations of a presentable $\D$ are of this form.
\end{rmk}

\subsection{Hom-categories}
\label{subsec:laxloc-homs}

\gutentag{00F7}%
For $X$, $Y \in \C$, the equivalence \eqref{eq:laxloc-adjunction} identifies
$R\colon \Hom_{\C}(X,Y) \to \Hom_{\D}(RX,RY)$
with
$\varepsilon_{X}^{\natural}\colon \Hom_{\C}(X,Y) \to \Hom_{\C}(LRX,Y)$,
naturally in $Y$, as
$R(k\varepsilon_{X}) \circ \eta_{RX} \simeq Rk$
by $\theta_{X}$. So $R$ is locally fully faithful if and only if every
$\varepsilon_{X}^{\natural}$
is fully faithful.

\begin{prop}
\label{prop:laxloc-homs}
\gutentag{00F8}%
Let $X \in \C$.
\begin{enumerate}
\item The transformation
$R\colon \Hom_{\C}(X,-) \Rightarrow \Hom_{\D}(RX,R-)$
of functors $\C \to \Cat$ has a right adjoint in $\func(\C,\Cat)$ if and only
if $\varepsilon_{X}$ has a left adjoint $l$, and the right adjoint is then
$l^{\natural}$, under
$\Hom_{\D}(RX,R-) \simeq \Hom_{\C}(LRX,-)$.
Its components are then coreflective if and only if
$l \dashv \varepsilon_{X}$
is a coreflection, e.g.\ if $R$ is locally fully faithful.
\item If $L \dashv R$ is a lax localization and $X = Ld$, then
$l = L\eta_{d}$
and
$l \dashv \varepsilon_{X}$
is a coreflection. That is, the functor
\begin{equation*}
\Hom_{\D}(d,RY) \longrightarrow \Hom_{\D}(Td,RY) \ , \qquad g \longmapsto
R\varepsilon_{Y} \circ Tg \ ,
\end{equation*}
is coreflective, in particular fully faithful, with right adjoint
$\eta_{d}^{\natural}$, naturally in $Y$.
\end{enumerate}
\end{prop}

\begin{proof}
\gutentag{00F9}%
(1) By \eqref{eq:laxloc-adjunction}, the transformation $R$ is
$\varepsilon_{X}^{\natural}$,
the image of $\varepsilon_{X}$ under the enriched
Yoneda embedding
$h\colon \C\opcat \to \func(\C,\Cat)$.
As in the proof of \cref{lem:laxloc-repr}, $h$ is fully faithful, takes
$l \dashv r$
to
$r^{\natural} \dashv l^{\natural}$
with the same unit and counit, and every adjunction between representables
arises this way; and a $2$-morphism of $\func(\C,\Cat)$ is invertible if and
only
if its components are. (2) By \cref{lem:laxloc-equiv}(4),
$L\eta_{d} \dashv \varepsilon_{Ld}$
is a coreflection, and under \eqref{eq:laxloc-adjunction},
$(L\eta_{d})^{\natural}$
becomes $\eta_{d}^{\natural}$ and $R$ becomes
$g \mapsto R(\varepsilon_{Y} \circ Lg) = R\varepsilon_{Y} \circ Tg$.
\end{proof}

\begin{exmp}
\label{ex:laxloc-homs}
\gutentag{00FA}%
Let $\D = \Cat$ and
$P = \{\emptyset \to \mathrm{pt}\}$,
so that
$\C := \Pc_{\emptyset,P}$
consists of the $\infty$-categories with an
initial object and the functors preserving it, and $R$ is the inclusion. As we
will see in \cref{ex:local-vs-pushout}(2), $R$ has the left adjoint
$L(d) = d^{\triangleleft}$,
and $L \dashv R$ is a normalized lax localization by \cref{cor:laxloc-P}.
\begin{enumerate}
\item For
$X = L(\mathrm{pt}) = \Delta^{1}$,
\cref{prop:laxloc-homs}(2) says
that the arrows $0 \to v$ out of an initial object form a coreflective full
subcategory of $\func(\Delta^{1},Y)$, with coreflector
$(u \to v) \mapsto (0 \to v)$.
For $Y = \Delta^{1}$ it is not reflective, as there is no map from
$(1 \to 1)$
to any
$(0 \to v)$.
\item Let $X$ be the walking retraction, with objects $0$, $1$ and morphisms
$i\colon 0 \to 1$,
$r\colon 1 \to 0$
with
$ri = \mathrm{id}_{0}$;
its object $0$ is initial. Then
$\varepsilon_{X}\colon X^{\triangleleft} \to X$
has no left adjoint $l$: for the cone point $\ast$,
$\mapp_{X}(1,\varepsilon_{X}(\ast)) = \mapp_{X}(1,0)$
is non-empty while only $\ast$ maps to $\ast$, so $l(1) = \ast$, and then $1$
would be initial. So
$R\colon \Hom_{\C}(X,-) \Rightarrow \func(X,-)$
has no right adjoint in $\func(\C,\Cat)$. Objectwise,
$\Hom_{\C}(X,Y) \simeq Y_{/0}$,
and $\func(X,Y)$ consists of the retractions
$(s,r)\colon A \rightleftarrows B$,
$rs \simeq \mathrm{id}_{A}$.
A morphism from
$(0 \rightleftarrows B',\ t\colon B' \to 0)$
to $(s,r)$ is a
$v\colon B' \to B$
with
$rv \simeq (0 \to A) \circ t$,
so $(s,r)$ has a coreflection if and only if the pullback of $r$ along
$0 \to A$
exists. Let $Y$ be the category generated by objects $0$, $A$, $B$, $C$ and
morphisms
$s\colon A \to B$,
$r\colon B \to A$,
$c\colon C \to 0$
and
$u,u'\colon C \to B$,
subject to $0$ being initial,
$rs = \mathrm{id}_{A}$
and
$ru = ru' = (0 \to A) \circ c$.
The morphisms to $0$ are $\mathrm{id}_{0}$ and $c$, the morphisms
$C \to B$ are $u$, $u'$ and
$(0 \to B) \circ c$,
and the endomorphisms of $C$ are $\mathrm{id}_{C}$ and
$(0 \to C) \circ c$.
So neither
$(0,\ 0 \to B)$
nor any $(C,v)$ is a pullback, and $(s,r)$ has no coreflection. It has no
reflection either, as there is no morphism $A \to 0$. So
$R\colon \Hom_{\C}(X,Y) \to \Hom_{\D}(RX,RY)$
is neither reflective nor coreflective.
\end{enumerate}
\end{exmp}

\begin{rmk}
\label{rmk:laxloc-continuous}
\gutentag{00FB}%
Following Johnstone--Joyal \cite{johnstonejoyal1982continuous}, who treat the
free completion under filtered colimits, call $X \in \C$ \emph{continuous} if
$\varepsilon_{X}$ has a left adjoint; for the ideal completion of posets, these
are the continuous dcpos. So \cref{prop:laxloc-homs}(1) says that
$R\colon \Hom_{\C}(X,-) \Rightarrow \Hom_{\D}(RX,R-)$
has a right adjoint if and only if $X$ is continuous. By
\cref{lem:laxloc-equiv}(4), every $Ld$ is continuous, and if $R$ is locally
fully faithful, every continuous $X$ is a retract of $LRX$, as
$l \dashv \varepsilon_{X}$
is then a coreflection. In \cref{ex:laxloc-homs}, $X$ is continuous if and only
if every object with a morphism to $0$ is initial, i.e.\ if and only if
$X \simeq Ld$
for some $d$, as we will see in \cref{ex:local-vs-pushout}(2); and retracts of
the $Ld$ are again of this form. So there continuity characterizes the
essential image of $L$. We do not know whether retracts of the $Ld$ are
continuous in general.
\end{rmk}

\begin{rmk}[the colax case]
\label{rmk:laxloc-colax}
\gutentag{00FC}%
Let $L \dashv R$ be a colax localization. Applying the results of this section
to the lax localization
$L^{\mathrm{co}} \dashv R^{\mathrm{co}}$
(\cref{subsec:laxloc-def}), and using that $(-)^{\mathrm{co}}$ exchanges
reflective and coreflective $1$-morphisms, and left Beck--Chevalley squares
with transposes of right Beck--Chevalley squares, gives the following.
We have
$\varepsilon L \dashv L\eta$
with counit $\varphi$ and
$\eta R \dashv R\varepsilon$
with unit $\theta^{-1}$, so $L$ makes the $\eta_{d}$ reflective and their
naturality squares $(g,Tg)$, drawn with $\eta_{d}$ horizontal, left
Beck--Chevalley; $\Sat_{L}$ is replaced by the $1$-morphisms that $L$ makes
reflective and the squares it makes left Beck--Chevalley. The modules form the
locally full sub-$(\infty,2)$-category $\D_{T}$ of the $x$ with $\eta_{x}$
coreflective and the $1$-morphisms whose transposed unit square is right
Beck--Chevalley (\cref{thm:kz-module-maps}), and
$\D_{T} = \Pc_{\Sigma_{\eta},\emptyset}$
in the notation that will be introduced in \cref{def:IP-local}:
the $i^{\natural}$
have fully faithful right adjoints. In (R1) and (R2), reflective $1$-morphisms
and their left adjoints are replaced by coreflective $1$-morphisms and their
right adjoints. The recognition criterion uses the lax limits
$(f \downarrow e)$
instead \cite[Rem.~3.3]{bourke2025recognition}. Finally,
$\Hom_{\C}(Ld,Y) \simeq \Hom_{\D}(d,RY)$
is a reflective full subcategory of
$\Hom_{\D}(Td,RY)$,
with reflector $\eta_{d}^{\natural}$.
\end{rmk}

